\chapter{Cryptographic Primitives}
\label{short:primitives}

Begin with the promise a cooperation problem needs. What should remain
private? What must be fixed before anyone learns more? What must another
party be able to check?

Then name the tool, the property it supplies, and the assumptions it asks us
to accept. Follow each use back to its definition and forward to the chapter
that depends upon it. A reader should be able to see where one promise ends
and another begins.

This map remains to be written in the companion. The research registry keeps
its place; the name of a primitive cannot fill an absent explanation.

\shortlimit{Choosing a cryptographic tool does not choose a worthwhile purpose.}
\companionref{3}{chapter.3}{An editorial placeholder for the map of primitives.}
